\chapter*{Profundización: cinco conexiones que preparan la investigación}
\addcontentsline{toc}{chapter}{Cinco conexiones para investigación}

\section*{1. ¿Qué se puede concluir de una cadena de Markov?}
\addcontentsline{toc}{section}{Cadenas de Markov y distribución estacionaria}
El modelo suavizado del ejemplo principal puede estudiarse como una cadena
homogénea de tres estados, con distribución inicial $\pi_0$ y matriz
\[
P=\begin{pmatrix}1/7&3/7&3/7\\1/2&1/3&1/6\\2/5&2/5&1/5\end{pmatrix}.
\]
Por construcción, $P_{ij}\ge0$ y $\sum_jP_{ij}=1$; cada fila define
una probabilidad. Como todas las entradas son estrictamente positivas,
la cadena es irreducible y aperiódica: desde cualquier estado se llega
a cualquier otro en un paso, y cada estado admite una transición a sí mismo.
Bajo estas condiciones existe una única distribución estacionaria $\pi$
con componentes estrictamente positivas y $\pi P=\pi$.

Para calcularla, escribimos $\pi=(u,v,w)$ y resolvemos
\[
\begin{cases}
 u=\tfrac17u+\tfrac12v+\tfrac25w,\\
 v=\tfrac37u+\tfrac13v+\tfrac25w,\\
 w=\tfrac37u+\tfrac16v+\tfrac15w,\\
 u+v+w=1.
\end{cases}
\]
Una de las primeras tres ecuaciones es redundante: las filas de $P$
suman uno. Tras sustituir y resolver el sistema lineal queda
\[
\boxed{\pi=\left(\frac{98}{281},\frac{108}{281},\frac{75}{281}\right).}
\]
Se verifica, por ejemplo, que
\[
\frac{98}{281}\frac17+\frac{108}{281}\frac12
+\frac{75}{281}\frac25
=\frac{14+54+30}{281}=\frac{98}{281}.
\]
El mismo procedimiento verifica las otras dos coordenadas.
Este cálculo revela la distinción entre \emph{probabilidad de transición}
y \emph{probabilidad de ocupación a largo plazo}: no son iguales por
definición. Una distribución inicial arbitraria evoluciona como
$\pi_t=\pi_0P^t$ y se aproxima a la estacionaria bajo las hipótesis
mencionadas. En lenguaje real, las probabilidades condicionadas suelen
cambiar con el contexto y no existe necesariamente una matriz homogénea
fija en ese sentido.

\begin{puente}
Este es el punto de contacto directo con el curso paralelo de Series
de Tiempo: cadenas de Markov, matrices estocásticas, distribución invariante
y pronóstico condicional. Más adelante podremos comparar esta representación
con modelos ocultos de Markov y modelos de estado, pero en S01 no es
necesario introducir sus algoritmos.
\end{puente}

\section*{2. Entropía condicional y tasa de información de una cadena}
\addcontentsline{toc}{section}{Entropía condicional y tasa de información}
La entropía conjunta de dos símbolos satisface
\[
H(X_1,X_2)=H(X_1)+H(X_2\mid X_1).
\]
Esto se demuestra al expandir
$-\sum_{i,j}p(i,j)\log_2 p(i,j)$ y usar
$p(i,j)=p(i)p(j\mid i)$:
\[
\begin{split}
H(X_1,X_2)&=-\sum_{i,j}p(i,j)\log_2p(i)\\
&\quad-\sum_{i,j}p(i,j)\log_2p(j\mid i)\\
&=H(X_1)+\sum_i p(i)H(X_2\mid X_1=i).
\end{split}
\]
El promedio $\sum_i p(i)H(X_2\mid X_1=i)$ es
la entropía condicional: mide cuánto falta por conocer sobre $X_2$
cuando se conoce el estado precedente.

Para la cadena estacionaria construida antes, su tasa de entropía es
\[
h(P,\pi)=-\sum_{i\in V}\pi_i\sum_{j\in V}P_{ij}\log_2P_{ij}.
\]
La fórmula surge porque el símbolo inicial contribuye un término finito
$H(X_1)$, mientras que cada transición, bajo estacionariedad, tiene
la misma entropía condicional media. En efecto,
\[
H(X_{1:T})=H(X_1)+\sum_{t=2}^T H(X_t\mid X_{t-1})
=H(\pi)+(T-1)h(P,\pi).
\]
Dividiendo por $T$ y dejando $T\to\infty$ se obtiene
$\lim_{T\to\infty}T^{-1}H(X_{1:T})=h(P,\pi)$.
Hay dos lecciones: el costo de representar el símbolo inicial
no desaparece en una cadena corta y la entropía marginal $H(\pi)$
no coincide necesariamente con la entropía condicional media.

\begin{ejemplo}[Un contraste con números sencillos]
Si las transiciones son independientes del estado anterior y cada fila
es $(1/2,1/2)$, entonces $H(X_t\mid X_{t-1})=1$ bit y la cadena
no aprovecha memoria. En cambio, para
$P=\left(\begin{smallmatrix}0.9&0.1\\0.1&0.9\end{smallmatrix}\right)$
la estacionaria es uniforme, $H(X_t)=1$ bit, pero
\[
H(X_t\mid X_{t-1})=-0.9\log_20.9-0.1\log_20.1\approx0.4690
\text{ bits}.
\]
El conocimiento del estado anterior reduce la incertidumbre de
la siguiente observación sin cambiar la distribución marginal.
\end{ejemplo}

\section*{3. Soporte, símbolos desconocidos y final de secuencia}
\addcontentsline{toc}{section}{Soporte y símbolos desconocidos}
El suavizado aditivo distribuye masa positiva entre los $K$ símbolos
incluidos en el vocabulario del modelo. \emph{No asigna automáticamente}
probabilidad a cualquier símbolo que no pertenezca a ese vocabulario.
En un modelo de palabras entrenado con $V$, una palabra nueva necesita
una decisión de representación: sustituirse por un símbolo
$\langle\mathrm{UNK}\rangle$, segmentarse en unidades menores o
admitirse mediante algún mecanismo adicional.

\begin{figure}[H]\centering
\begin{tikzpicture}[>=Stealth,font=\small,
 box/.style={draw=azul,fill=azulclaro,rounded corners,align=center,text width=3.1cm,minimum height=9mm}]
\node[box](a)at(0,0){Palabra de prueba\\\texttt{revestidor}};
\node[box](b)at(4.1,1){Vocabulario de palabras\\\texttt{<UNK>}};
\node[box](c)at(4.1,-1){Vocabulario subléxico\\varias unidades};
\draw[->,thick] (a)--(b);\draw[->,thick] (a)--(c);
\end{tikzpicture}
\caption{Dos políticas distintas para un símbolo no visto; no son intercambiables.}
\end{figure}

Si además la longitud de la cadena es variable, necesitamos fijar
cómo termina la secuencia. Una construcción usual introduce un símbolo
de fin $\langle/\mathrm{s}\rangle$ y asigna probabilidades a las
transiciones hasta encontrarlo. De este modo, bajo hipótesis de
terminación adecuadas, el modelo determina una ley sobre cadenas
finitas de longitudes variables. Si se omite la probabilidad de parar,
la suma de probabilidades de cadenas completas puede no estar
bien definida.

\begin{ejemplo}[Parada con probabilidad geométrica]
Considérese un generador que emite un símbolo ordinario con probabilidad
$1-r$ y termina con probabilidad $r\in(0,1)$ en cada paso, con decisiones
independientes. Si $N$ es el número de símbolos ordinarios previos al fin,
\[
P(N=n)=(1-r)^nr,\quad n=0,1,2,\ldots,
\qquad \mathbb E[N]=\frac{1-r}{r}.
\]
Al limitar la generación a un máximo de $m$ símbolos mediante un arnés,
la probabilidad de truncamiento es $P(N\ge m)=(1-r)^m$.
Esta cuenta es una idealización explícita: en un LLM real la probabilidad
de fin depende de la historia y la política de decodificación.
\end{ejemplo}

\section*{4. Códigos ideales y límites de interpretación}
\addcontentsline{toc}{section}{Código ideal y pérdida logarítmica}
La cantidad $-\log_2q(x)$ se interpreta como longitud de código ideal
según $q$. Esto motiva tanto la entropía como la divergencia KL:
si los datos siguen $p$ pero se utiliza $q$,
\[
\mathbb E_p[-\log_2q(X)]-\mathbb E_p[-\log_2p(X)]
=D_{\rm KL}(p\|q).
\]
Es el exceso esperado de longitud ideal, no una promesa de que
cualquier compresor real tenga exactamente esa longitud en un archivo.
Los códigos binarios de longitud entera deben cumplir restricciones
como la desigualdad de Kraft y agregan sobrecostos de cabecera,
fronteras y representación.

\begin{ejemplo}[Cuando la entropía no es el costo de un símbolo concreto]
Para $p=(1/2,1/4,1/4)$, el símbolo de probabilidad $1/2$ contribuye
$1$ bit si ocurre, mientras que cada símbolo de probabilidad $1/4$
contribuye $2$ bits. La entropía $1.5$ bits es un \emph{promedio esperado},
no la longitud de cada realización. Con un código de prefijo
$\{0,10,11\}$ las longitudes individuales son exactamente $(1,2,2)$,
y el promedio coincide con 1.5 bits.
\end{ejemplo}

La pérdida logarítmica es una regla de puntuación propia: su riesgo
esperado se minimiza al informar la distribución verdadera.
La demostración ya está contenida en la identidad de entropía cruzada
más Gibbs. Por ello tiene sentido usarla para evaluar predicciones
probabilísticas, pero no sustituye métricas de exactitud, fidelidad o
utilidad operacional cuando esas son las preguntas científicas.

\section*{5. Incertidumbre experimental y unidad de remuestreo}
\addcontentsline{toc}{section}{Unidad estadística de comparación}
Supongamos que se quiere comparar dos procedimientos A y B en $m$
documentos reservados. Si $L_{A,j}$ y $L_{B,j}$ denotan las pérdidas
en el mismo documento $j$, definimos
\[
d_j=L_{A,j}-L_{B,j},\qquad \overline d=\frac1m\sum_{j=1}^m d_j.
\]
Un análisis pareado retiene la estructura de comparación: cada documento
aporta una diferencia. Si las unidades documentales son razonablemente
independientes, su error estándar clásico es $s_d/\sqrt m$, donde
\[
s_d^2=\frac1{m-1}\sum_{j=1}^m(d_j-\overline d)^2.
\]
Este error estándar y los intervalos que se construyen con él no
resultan válidos por el mero hecho de tener millones de tokens: los
símbolos dentro de un mismo documento comparten dependencias.

En investigación con textos técnicos, varios documentos del mismo pozo,
incidente o empresa también pueden estar correlacionados. En ese caso la
unidad de remuestreo puede necesitar elevarse hasta el pozo, incidente o
fuente, según la pregunta y el proceso de muestreo. Una comparación que
ignora esto puede aparentar una precisión mayor de la real.

\begin{ejemplo}[Diferencias pareadas construidas para ilustrar el método]
Considérense cuatro documentos con diferencias en BPC
$(-0.20,0.10,-0.10,0.00)$, calculadas como A menos B.
Su promedio es $-0.05$ BPC. Los cuadrados respecto de la media son
$0.0225,0.0225,0.0025,0.0025$; la suma es $0.05$.
La varianza muestral es $0.05/3\approx0.01667$ y el error estándar
$\sqrt{0.01667/4}\approx0.06455$. Con solo cuatro documentos,
el valor puntual no debe venderse como efecto establecido. El ejemplo es
\textbf{ILUSTRATIVO}, no procede de una prueba real.
\end{ejemplo}

\section*{6. Cuidado con una supuesta ventaja metodológica}
\addcontentsline{toc}{section}{Separar hipótesis de evidencia}
Un tokenizador puede producir menos unidades por palabra sin que
su modelo probabilístico correspondiente alcance menor pérdida de
predicción; la arquitectura, la cantidad de datos y el tratamiento
de caracteres pueden alterar esa relación. Por eso conviene separar
métricas \emph{descriptivas de segmentación} de métricas \emph{predictivas}.
Si comparamos dos modelos completos, podremos describir diferencias
de tasa; si queremos atribuir la diferencia a BPE, habrá que diseñar un
experimento en el que los otros factores se mantengan comparables.

Esta cautela también explica por qué la matriz IC10 es una \emph{semilla}
de investigación, no un resultado preestablecido. Un resultado nulo,
negativo o contrario a la hipótesis puede ser científicamente útil si
la medición fue defendible. El paso siguiente, en S02, será revisar
candidatos y disponibilidad de datos antes de elegir un problema.
